\documentclass[conference]{IEEEtran}
\usepackage[utf8]{inputenc}
\usepackage[T1]{fontenc}
\usepackage{cite}
\usepackage{amsmath,amssymb}
\usepackage{graphicx}
\usepackage{booktabs}
\usepackage{url}

\title{Isolated Downstream Validation and Single‑Repair Effects in Intent‑to‑Configuration NetOps: A Preregistered Paired Study}

\author{
\IEEEauthorblockN{Autonomous AI Researcher}
\IEEEauthorblockA{CNSM 2026 Agentic AI Researcher Track}
}

\begin{document}

\maketitle

\begin{abstract}
We report a preregistered paired confirmatory experiment (study cand-003-stud-001) that isolates purely downstream deterministic validation plus a bounded single-repair intervention under shared\_initial\_candidate semantics. Using 1,000 deterministic intents (1,000 paired episodes) drawn from a synthetic workflow stress corpus and a deterministic digital‑twin validator (deterministic\_netops\_validator\_v5), exactly one sampled initial generation per task was deterministically reused by both arms; any paired outcome difference can therefore arise only from the validator decision and at most one permitted repair call (guarded arm). Observed outcomes: baseline safe‑deploy = 995/1000 = 0.995; guarded safe‑deploy = 1000/1000 = 1.000; discordants n10 = 5 (guarded only), n01 = 0 (baseline only); preregistered primary test exact two‑sided McNemar p = 0.0625 (not < 0.05). Preregistered paired bootstrap percentile 95\% CI (10,000 resamples, seed = 1729) = [0.001, 0.01]. Total hosted model calls recorded = 1005 (1000 shared initial + 5 repair). We present the exact McNemar as the confirmatory inference, the bootstrap interval as a preregistered descriptive summary, and interpret the small absolute guarded benefit in light of sparse discordant counts, ceiling effects, validator scope, and the shared\_initial\_candidate design.
\end{abstract}

\section{Introduction}
Automating translation from operator intent to device configuration promises operational efficiency but raises safety risks: incorrect command ordering, transient invariant violations, and unauthorized touches of protected fields can convert plausible plans into damaging live changes. Prior empirical work documents frequent LLM hallucination, prompt sensitivity, and structured‑output failures that reduce correctness in code- and config-like generation and motivates verifier‑augmented architectures that detect and sometimes repair unsafe outputs [1],[2].

This paper reports a preregistered paired experiment explicitly designed to isolate the incremental effect of a deterministic validator and a single bounded repair invocation when both experimental arms reuse the exact same initial sampled candidate (generation\_semantics = shared\_initial\_candidate). Our central empirical estimand is the paired success‑rate difference (guarded - baseline) computed on complete pairs. Under shared\_initial\_candidate semantics no first‑pass sampling differences exist between arms; therefore any paired outcome difference can only arise from the deterministic validator decision and the allowed at‑most‑one repair call. This causal framing differs from generate\ensuremath{\rightarrow}validate\ensuremath{\rightarrow}repair evaluations that mix validator effects with independent sampling or decoding changes and makes the estimand precisely: baseline single‑shot/initial validity versus final guarded‑pipeline validity after deterministic validation and permitted repair.

Contributions. (1) A preregistered paired confirmatory study (Npairs = 1,000) with deterministic task selection from a challenging synthetic NetOps workflow stress corpus and archived validator traces. (2) Full execution accounting and explicit runtime configuration reporting (declared model = gpt‑5‑mini; model configuration recorded temperature = null/unset --- null/unset is not evidence of deterministic hosted‑model sampling and does not imply temperature = 0). (3) Transparent reporting of the preregistered primary inference (exact two‑sided McNemar, p = 0.0625), the preregistered paired bootstrap descriptive summary (95\% percentile CI = [0.001, 0.01], 10,000 resamples, seed = 1729), and an evidence‑grounded interpretation that accounts for sparse discordant counts, near‑unity baseline, and validator scope.

\section{Related Work}
Empirical studies of LLMs on structured generation consistently report sensitivity to prompting and decoding choices and a nontrivial rate of syntactic and semantic errors that affect downstream executability [3],[4]. Several lines of work address this by (a) enforcing constraints at decode time (constrained/schema decoding) and (b) post‑generation verification and repair. Evidence from constrained‑decoding method papers shows improvements in executable outputs in structured domains (for example, cheminformatics), though that evidence is domain‑specific and depends critically on the quality of the constraint tooling [11],[52].

Separate generate\ensuremath{\rightarrow}validate\ensuremath{\rightarrow}repair architectures and digital‑twin validators have been proposed and evaluated in NetOps and related settings; these works provide system prototypes and benchmark results showing that verifier-augmented pipelines can reduce observable errors under controlled scenarios, but measured gains vary with validator coverage and benchmark design [5],[6],[7]. Importantly, many prior evaluations conflate validator/repair effects with independent changes in generation (distinct prompts or decoding strategies), which complicates causal attribution. Method papers therefore motivate an isolating experimental design that holds generation constant while measuring downstream validator/repair marginal returns [5],[7].

Agentic and multi‑stage NetOps architectures emphasize governance, auditability, and protocol‑mediated execution (proof‑carrying traces, canaries, rollbacks) as complements to any single model's behavior; those works argue that operational reliability depends on validator fidelity and orchestration policies rather than on the LLM alone [8],[9],[10]. Our study complements these prior lines by holding initial generation constant (shared\_initial\_candidate) and measuring only the deterministic downstream effect of a validator plus a bounded single repair invocation, thereby answering a narrowly specified causal question recommended by the literature: what marginal safety return does a deterministic digital‑twin validator and one permitted repair provide when the generator output is fixed?

\section{Methodology}
Preregistration and confirmatory design. The preregistration (study id cand‑003‑stud‑001) specified a paired binary design with task\_count = 1,000 complete pairs, generation\_semantics = shared\_initial\_candidate (exactly one sampled initial generation per task deterministically reused by both arms), initial\_generation\_calls\_per\_task = 1, maximum\_repair\_calls\_per\_task = 1, primary estimand = paired\_success\_rate\_difference\_guarded\_minus\_baseline, and primary confirmatory test = exact two‑sided McNemar. The paired bootstrap percentile 95\% CI (10,000 resamples, seed = 1729) was prespecified as a descriptive summary. The preregistration explicitly set independent\_condition\_generation = false; no independent first‑pass generations for alternate prompting or decoding were executed. These design choices were frozen before any hosted‑model calls.

Task corpus, deterministic selection, and difficulty taxonomy. Tasks come from the synthetic challenging\_workflow\_stress\_v1 corpus and were deterministically selected to produce 1,000 unique intents. Each task manifest includes initial\_state, expected\_final\_state, required\_sequence, allowed\_touched\_fields, and a difficulty annotation (levels include 'very\_hard' and 'extreme'). Representative workflow patterns include safe\_access\_service\_transfer, failover\_then\_offline\_reconfiguration, rolling\_access\_vlan\_migration, and composed\_redundancy\_and\_access\_maintenance; these patterns expose ordering and transient invariants that the validator checks. Selection was deterministic (see preregistration sampling rule) so the analysis treats the resulting 1,000 pairs as the preregistered confirmatory set.

Adapter, model configuration, and validator semantics. The adapter family was hosted\_netops\_gvr\_v1 and the declared provider/model was gpt‑5‑mini (provider: openai\_responses). The recorded model\_configuration shows temperature = null/unset; per preregistered reporting we explicitly state that null/unset is not evidence of deterministic hosted‑model sampling and does not imply temperature = 0. Exactly one sampled initial generation per task was deterministically reused by both arms (shared\_initial\_candidate semantics); the guarded arm applies deterministic validator checks and may invoke at most one repair call if the validator rejects the initial candidate. The validator deterministic\_netops\_validator\_v5 is a digital‑twin emulator that parses candidate command sequences and checks required command presence/order, protected‑field noninterference, and workflow invariants; it emits a binary safe‑deploy decision and structured traces (validation\_before, repair\_applied, validation\_after).

Causal interpretation under shared\_initial\_candidate. Under shared\_initial\_candidate semantics the baseline outcome is the initial single‑shot validity of the shared candidate and the guarded outcome is the final validity after deterministic validation and any permitted single repair. Because the pair reuses the same initial candidate, any observed discordance (n10 or n01) can only arise from validator-triggered repair (or validator decision differences applied deterministically). The preregistered design therefore isolates downstream validator/repair marginal effects but cannot be used to evaluate prompt or decoder interventions; planned independent-generation arms were preregistered but not executed in this run.

\section{Results}
Execution accounting. Planned episodes = 2,000; completed = 2,000; complete\_pairs = 1,000; failed episodes = 0. Provider call audit shows hosted model call event count = 1,005 consisting of 1,000 shared initial generations and 5 repair invocations (stage counts recorded in the deterministic reconciliation artifact).

Primary paired outcomes (preregistered). The preregistered paired contingency table is: n11 = 995 (both succeeded), n10 = 5 (guarded succeeded, baseline failed), n01 = 0 (baseline succeeded, guarded failed), n00 = 0 (both failed). Observed success rates: baseline = 995/1000 = 0.995; guarded = 1000/1000 = 1.000; paired difference (guarded - baseline) = 0.005. The exact two‑sided McNemar p = 0.0625 (primary prespecified confirmatory test; \ensuremath{\alpha} = 0.05: not significant).

Paired bootstrap descriptive summary. The preregistered paired bootstrap percentile 95\% CI (10,000 resamples, seed = 1729) computed by the registered executor is [0.001, 0.01]. This prespecified descriptive interval excludes zero in this run despite the confirmatory p‑value exceeding 0.05; this discrepancy is a known consequence of discreteness when discordant counts are extremely sparse.

Repair mechanics and representative diagnostics. Exactly five validator‑triggered repairs occurred and produced the five discordant improvements. Archived validator traces for repaired pairs document rejection of the shared initial candidate for rule‑level violations (notably ordering violations in required\_sequence and unauthorized touches of protected interfaces), a single repair call that produced a repaired candidate, and subsequent revalidation that accepted the repaired candidate. The repaired episodes are concentrated among higher difficulty patterns annotated as 'very\_hard' or 'extreme' and involve multi‑step required\_sequences. Total hosted model calls = 1,005 (1,000 shared initial + 5 repair) and all deterministic reconciliation checks passed.

Compact execution/accounting table (summary). Planned episodes = 2,000; completed = 2,000; complete pairs = 1,000; model calls used = 1,005; repair calls invoked = 5; baseline successes = 995; guarded successes = 1,000; paired difference = 0.005; exact two‑sided McNemar p = 0.0625; paired bootstrap 95\% CI = [0.001, 0.01] (10,000 resamples, seed = 1729).

\section{Discussion}
Confirmatory vs descriptive evidence and discrete inference nuance. The preregistered decision rule treats the exact two‑sided McNemar (p = 0.0625) as the confirmatory outcome and, under the prespecified two‑sided \ensuremath{\alpha} = 0.05, the study does not claim confirmatory support for a nonzero paired increase. The preregistered bootstrap CI ([0.001, 0.01]) is a descriptive summary produced by the registered executor and is reported to aid interpretation; it does not override the preregistered confirmatory conclusion. When discordant counts are sparse (here n10 + n01 = 5), the null distribution of exact paired tests is concentrated on a few small symmetric outcomes and two‑sided exact p‑values can remain above 0.05 while resampling‑based percentile intervals exclude zero. This is a sampling‑distribution phenomenon, not a contradiction in the underlying recorded counts.

Power, ceiling effects, and sensitivity. The preregistered power calculation assumed baseline\ensuremath{\approx}0.60 and targeted an absolute increase of 0.15; under those assumptions task\_count = 1,000 provided >80\% power for the McNemar-style test. Observed baseline = 0.995 greatly departs from that assumption and concentrates mass in the concordant success cell (n11 = 995), leaving only five discordant pairs for inference. Detecting the observed tiny absolute increase (0.005) against a near‑unity baseline requires either (a) dramatically larger paired sample sizes or (b) design choices that increase discordant prevalence (for example, independent generation arms, injected failure cases, or a larger allowed repair budget). We therefore interpret the null confirmatory result in the context of a mismatch between preregistered operating assumptions and the sampled operating point.

Failure‑mode diagnostics and difficulty stratification. Archived validator traces show repaired cases shared rule‑level failure modes: command ordering violations in required\_sequence and unauthorized touches of protected interfaces. Repaired tasks are drawn from complex multi‑step workflows (patterns cataloged as safe\_access\_service\_transfer, failover\_then\_offline\_reconfiguration, rolling\_access\_vlan\_migration, and composed\_redundancy\_and\_access\_maintenance) and were annotated 'very\_hard' or 'extreme'. This qualitative pattern---repairs concentrated in high‑difficulty, multi‑step tasks---is consistent with the expectation that dependency constraints raise ordering vulnerabilities. However, with only five repaired cases we cannot robustly estimate per-pattern repair efficacy; the evidence supports only the narrower, qualified statement that deterministic validator+single repair averted a small number of unsafe initial candidates concentrated in higher‑difficulty workflows.

Design tradeoffs and scope. The shared\_initial\_candidate paired design intentionally isolates downstream validator and single‑repair effects and avoids conflation with prompt/decoding or sampling interventions; this yields a precise causal estimand (baseline initial validity versus final guarded validity after deterministic validation and up to one repair) but precludes measuring generator‑level interventions. The preregistered plan included independent‑generation arms and repair‑budget sweeps as follow‑ups; those arms were not part of the frozen execution and therefore are not evaluated here. The deterministic digital‑twin validator provides reproducible rule‑based adjudication but its coverage determines construct validity: omitted semantic checks would inflate measured success.

Operational implications and concrete next steps. Artifact‑grounded evidence shows that a deterministic validator with a bounded single repair call can avert some unsafe outcomes in synthetic NetOps workflows, but measurable marginal gains are small when baseline outputs are already highly valid and when the repair budget is highly constrained. Supported, concrete follow‑ups (all preregistered or explicitly proposed in the archived manifest) include: (1) executing independent first‑pass generation arms to measure generator+validator interactions (preregistered but not run here); (2) sweeping repair budgets (k > 1 repairs) to quantify marginal returns per additional repair call; (3) expanding validator rule suites, injecting targeted failure cases, and varying topology heterogeneity to probe external validity; and (4) evaluating multiple model families and prompt/configuration variants. Each follow‑up should prespecify baseline failure ranges and adjust power calculations to the expected operating regime.

Threats to validity and limitations (summary). Construct validity depends on the validator rule set; omitted semantic checks reduce measured failures. Internal validity: shared\_initial\_candidate semantics correctly isolate downstream effects but prevent evaluating generator interventions. External validity: synthetic corpus and single declared model (gpt‑5‑mini) limit generalization to heterogeneous production environments. Conclusion validity: extreme baseline success produced sparse discordant counts and discrete p‑value behavior, reducing sensitivity for small absolute effects. These limitations were prespecified and are documented in the archived manifest.

\section{Limitations}
\begin{itemize}
\item Single‑model scope: only the declared model gpt‑5‑mini was evaluated; results do not generalize across model families or versions.
\item Synthetic corpus and deterministic validator: tasks come from a synthetic stress corpus and validation used a deterministic digital twin (deterministic\_netops\_validator\_v5); external validity to heterogeneous vendor/legacy environments is limited.
\item Shared\_initial\_candidate semantics: the paired comparison isolates downstream validator/repair effects only; it does not measure differences that would arise from independent first‑pass generations or alternative prompting arms that were not executed (those arms were not part of the frozen preregistration).
\item Low baseline failure prevalence: baseline success = 99.5\%; the preregistered power plan assumed baseline\ensuremath{\approx}0.60 and a target effect of 0.15, so the observed baseline greatly reduces sensitivity to detect small absolute gains. This mismatch is reported explicitly and reduces the study's ability to detect very small practical effects using the preregistered confirmatory test.
\item Sparse repairs and interpretation limits: only five validator repairs were invoked, constraining inference about repair efficacy across failure modes.
\item Statistical nuance: the preregistered primary test is exact two‑sided McNemar (p = 0.0625). The paired bootstrap percentile CI [0.001, 0.01] is a preregistered descriptive summary but is sensitive to sparse discordant counts; a one‑sided McNemar (p = 0.03125) was computed post‑hoc and is explicitly labeled exploratory because it was not prespecified in the frozen preregistration.
\item Construct validity: the safe‑deploy binary depends on the emulator/validator's rule set; semantic and emergent multi‑device policy violations may be underrepresented in the validator's modeled constraints.
\end{itemize}

\section{Conclusion}
In a preregistered paired study with shared\_initial\_candidate semantics on a 1,000‑task synthetic NetOps stress corpus, a guarded generate\ensuremath{\rightarrow}validate\ensuremath{\rightarrow}repair pipeline produced a small absolute paired improvement (0.005) relative to a baseline single‑shot generator; the preregistered exact two‑sided McNemar did not meet two‑sided \ensuremath{\alpha} = 0.05 (p = 0.0625). Only five validator repairs were invoked and total hosted model calls recorded = 1,005. The preregistered confirmatory result is the exact McNemar outcome; the paired bootstrap CI and one‑sided test are reported as descriptive and exploratory respectively. Artifact‑grounded evidence indicates downstream validator/repair safeguards can prevent some unsafe outcomes in structured NetOps workflows, but measurable marginal gains are constrained when baseline outputs are already highly valid and when the paired design isolates downstream effects. Follow‑ups that vary generator semantics, repair budgets, validator coverage, and model families are necessary to map marginal returns in realistic operating regimes.

\section*{Disclosure Statement}
Declared model and provider: gpt‑5‑mini (provider: openai\_responses). Orchestration and validator: hosted\_netops\_gvr\_v1 with deterministic\_netops\_validator\_v5 as the validator/digital twin. Preregistration: study cand‑003‑stud‑001 (frozen prior to any model calls and recorded in the archived execution manifest). Immutable initial master prompt reference archived at provenance/master\_prompt.txt (SHA‑256 = 1872df1e\allowbreak{}1805d2d9\allowbreak{}6940456c\allowbreak{}a016bd66\allowbreak{}5d1d5196\allowbreak{}add77f5a\allowbreak{}cdf1582b\allowbreak{}b39b15ba). Model configuration recorded temperature = null/unset (null/unset does not imply deterministic sampling or temperature = 0). Human role and autonomy boundary: after the autonomy lock the research run executed the preregistered experiment, analysis, and manuscript drafting; no human manually edited the manuscript technical content after the autonomy lock.

\begin{thebibliography}{99}
\bibitem{ref1} Dang Anh-Hoang, Vu Tran, Le-Minh Nguyen, "Survey and analysis of hallucinations in large language models: attribution to prompting strategies or model behavior", 2025, 10.3389/frai.2025.1622292
\bibitem{ref2} Ziyao Zhang, Chong Wang, Yanlin Wang, Ensheng Shi, Yuchi Ma, Wanjun Zhong, Jiachi Chen, Mingzhi Mao, Zibin Zheng, "LLM Hallucinations in Practical Code Generation: Phenomena, Mechanism, and Mitigation", 2025, 10.1145/3728894
\bibitem{ref3} Sakshyam Banjade, "Schema-Constrained LLM Planning for Executable Molecular Workflows: An Intent-to-Execution Infrastructure for Cheminformatics", 2026, 10.26434/chemrxiv.15005091/v1
\bibitem{ref4} Ari Harrison, "Protocol-Mediated Execution Governance for LLM Agent Systems", 2026, 10.2139/ssrn.6426760
\bibitem{ref5} https://arxiv.org/abs/2604.18233
\bibitem{ref6} Rufat Asadli, Benjamin Hoffman, Ioannis Protogeros, Laurent Vanbever, "Evaluating Agentic Configuration Repair for Computer Networks", 2026
\end{thebibliography}

\end{document}
